% % % ----------------------
% % \subsection{Gospel features and physical constraints}
% % % ----------------------

	% % We'll now turn our attention to the relationship between features of the gospel and physical constraints.
	% % Our goal is to better see why naturalistic explanations are important in thinking about religious phenomena.
	% % The philosophy will get a bit heavier here,
		% % as we make our way deeper into understanding how to think about many issues in the LDS church.
		
	% % -----
	% \subsubsection{Constraints on design and purpose}
	% % -----
	
		% Let's go back to an example I used in an earlier essay,
			% about walking through a half-finished house.
		% Because you know the designs and purposes of the builders,
			% you are able to grasp why there's a hole in the center of the room---%
				% it's going to be a kitchen island.
		% You understand many other features in the same way,
			% by thinking about the reasons the builders had in mind when they did certain things.
			
		% Now let's imagine that we walk into a friend's house just after the builders finish.
		% Our friend has been telling us to come over for a while because she can't wait to show us her new kitchen.
		% We walk into the house,
			% go through the family room,
				% and into the kitchen.
		% It's a huge room with lots of cabinets and counter space.
		% The dining area is just past the last row of cabinets,
			% and our friend has installed a big table she can't wait to use.
		% But then we notice something weird.
		% Just to the right of the table,
			% the wall juts out in a big semicircle.
		% The protrusion encroaches into the dining area and takes up some space that could have been more kitchen.
		% It doesn't have any windows.
		% The semicircle is an odd feature of the house---%
			% you've never seen anything like it,
				% and it doesn't seem like it has much purpose.
				
		% You decide to reflect for a moment before asking your friend why it's there.
		% In your mind you start to come up with reasons why someone might have such a feature in their house.
		% Your friend is a bit eccentric,
			% so it could just be her unusual tastes in floor plans.
		% You also know that the company which built her home has a reputation for innovative designs,
			% so maybe they put it there to start a trend.
		% On the other hand,
			% maybe the semicircle \textit{does} have a purpose---%
				% there might be a door leading inside to a pantry or to another storage area.
		% Now that you're thinking about it,
			% you realize it could have many more purposes than you first realized.
		% It could be there for lots of reasons.
		
		% So you decide to ask your friend what the semicircle is all about.
		% ``Come outside and I'll show you,''
			% she says.
		% She takes you out a side door to the rear of the house,
			% to a spot just outside where you were standing in the kitchen before.
		% There in front of you is an enormous boulder,
			% in the rough shape of a semicircle,
				% sticking out from a hill.
				
		% Our friend explains that she wanted to build the house on this property,
			% at this exact spot,
				% and she wanted the kitchen to be up against the hill.
		% But the boulder was in the way.
		% They tried dynamite,
			% a wrecking ball,
				% and gentle massages,
					% but nothing worked---%
						% the rock was stuck in place.
		% No matter what they did,
			% the excavators could not get rid of the boulder.
		% They had to give up.
		% The excavators told the home designer,
			% and they told our friend.
		% They gave her a choice:
			% either change your plans and put the kitchen or even the whole house somewhere else,
				% or built it around the boulder.
		% Our friend was adamant about the location of both the house and the kitchen,
			% so the designers and builders did what she said.
		% They put the kitchen where she wanted it and just built around the boulder.
		% They made a semi-circular indentation into the frame of the house to accommodate the rock,
			% and now it juts into the kitchen.
		
		% Once our friend tells us this story,
			% our view of the semicircle inside the house is going to change.
		% When we first saw it we assumed it had a more expicit purpose,
			% like storage or an avant-garde statement to people who like square kitchens.
		% But nothing like that turned out to be true---%
			% the semicircle was there because it \textit{had} to be.
		% If the kitchen was going to be right there,
			% it couldn't \textit{not} turn out like that.
			
		% Now,
			% in a sense,
				% the semicircle does have a purpose:
					% accommodating the boulder,
						% of course!
		% But this purpose is very different from that of the kitchen island's location,
			% for example.
		% Our friend could have put the island anywhere---%
			% she could have chosen any spot in the kitchen or even out of it.
		% She chose a certain spot for certain purposes,
			% like being close to the sink and the fridge and the oven at the same time.
		% The location of the island has that purpose,
			% and construction was carried out with that design in mind.
			
		% Compare that design to the design of the semicircle.
		% Sure,
			% it was designed to be that way;
				% sure,
					% it had the purpose of making room for the boulder so the kitchen could go against the hill.
		% But in reality,
			% as long as our friend wanted the kitchen against the hill,
				% the semicircle \textit{had} to be there.
		% She couldn't have moved it around the house like she could the island,
			% choosing its location on the basis of various reasons.
		% Had she wanted to eliminate the semicircle,
			% she would have had to change the entire design of the house or move the whole thing to a new place.
		% Changing the island's location didn't require anything like that.
		
		% We can think of the boulder as a \textit{physical constraint} on the design of the house.
		% The rock was a given;
			% it could not change;
				% they had to build the house around it.
		% Although the semicircle in the house did have a purpose involving the rock,
			% the semicircle in fact could not have been different.
		% It wasn't a design choice like where to put the island or how high the oven should rise above the ground.
		% Because of the physical constraint of the rock,
			% the semicircle could not have been otherwise.
		% The house had to end up built like that.			
			
	% -----
	\subsubsection{Strong and weak senses of design}
	% -----
	
		% Now we're considering two senses of the word ``design'' (or ``purpose,''
			% ``intention,''
				% and so on).
		% In the sense that describes the location of the kitchen island,
			% the reasons for the design don't have anything to do with physical constraints like the rock.%
				% \footnote{Basic physical constraints are still acting, of course. The island had to be located in space and time, it had to be made of matter, and so on. We're ignoring those very basic constraints...for now.}
		% The reasons are all about convenience,
			% ease of cleanup,
				% and so on.
		% This opens up many choices of location for the island---%
			% the design is unfettered by most physical limits,
				% such as whether there's a mountain in the way or not.
				
		% We'll call this \textit{design in the strong sense}.
		% We get design in the strong sense when the reasons explaining a design choice aren't dictated by physical constraints,
			% or when the design choice is a real choice from among many distinct and viable options.
		% The location of the island is designed in this strong sense.
		% It could have been anywhere,
			% and our friend wasn't limited by the presence of a boulder or lake in making her choice.
			
		% This situation is different from design in the weak sense,
			% which describes the ``design'' of the semicircle.
		% There's no doubt the semicircle has a purpose,
			% but like we said before,
				% it had to be that way.
		% Other than moving the kitchen or the whole house,
			% there were no options.
		% And since the location of the kitchen and the house were not negotiable,
			% the physical constraint dictated the purpose of the semicircle.
		% There was design, 
			% but it wasn't unfettered.
		% Weak design involves more salient physical constraints.
			
		% Another way to see the distinction between the strong and weak senses of design is to look at how they work in explaining things.
		% When we ask our friend why the island is where it is,
			% she'll say something like,
				% ``Because it was the most convenient spot'' or ``I just wanted it close to everything.''
		% Those reasons of hers factored into the island's location.
		% That's the strong sense of design.
		% But when we ask why the semicircle is there,
			% she has to say,
				% ``It had to be, because of the boulder.''
		% The explanation for the semicircle refers to the physical constraint which limited its design.
		% And that's design in the weak sense.
		
		% The distinction between strong and weak design turns out to be really useful.
		% It's true that the semicircle has a design,
			% and there are reasons for it.
		% But when we ask for the reasons we get something surprising---%
			% some more basic and unalterable fact about the background or environment,
				% rather than a reason relating to human-comprehensible purposes or desires.
		% That's a very different kind of explanation from one that just involves our friend choosing from a bunch of possible locations with no constraints on choice.
		% Separating the strong and weak senses of design helps us both respect this difference between explanations,
			% and explain it.
			
		% Let's see one area in gospel thought where we can apply the distinction between strong and weak design.
		% We'll use another example we've followed before too,
			% based on the observation that God made only men bearers of the priesthood.
		% God's action is that he bestowed the priesthood this way,
			% which means he ``designed'' things to be that way,
				% or that he ``purposed'' it to happen like that.
		% In the vaguest way,
			% it's true that God designed for only men to bear the priesthood.
		% But we now know that there are multiple senses of ``design,''
			% so our statement about God's ``design'' is ambiguous until we specify which sense we mean.
		% We can't know what to make of this ``design'' until we understand how strong it is.
		% Let's take a look at some possibilities:
		
		\begin{table}[H]
		\centering
		\begin{tabular}{p{1cm}p{4cm}p{9cm}p{1cm}}
		\textbf{Sense} & \textbf{Statement} & \textbf{Interpretation / Explanation} & \textbf{Design?} \\ \hline
		\#1 & Men bear the priesthood, and not women. & God could have given the priesthood to anyone, including women, but he willed to give it to men only, for either no reason at all or for reasons we can never know or understand. & Strongest \\ \hline
		\#2 & Men bear the priesthood, and not women. & God could have given the priesthood to anyone, including women, but he chose to gave it to men instead, knowing them to be more capable/required to serve others/etc. & Stronger \\ \hline
		\#3 & Men bear the priesthood, and not women. & God looked down the road into eternity and recognized that only men are able to bear the priesthood; he therefore could \textit{not} have given it to anyone, and therefore made men the bearers of the priesthood, and not women. & Weak \\ \hline
		\#4 & Men bear the priesthood, and not women. & God would have given the priesthood to everyone. He looked at human society in 1829, however, and saw that men were in charge of most things, and that many people would neither accept nor even consider legitimate a religion which gave priesthood and leadership roles to women. Knowing the needs of the early church, he therefore could \textit{not} have given it to anyone, at least at that time, and so made men the bearers of the priesthood. & Weak \\
		\end{tabular}
		% \caption{\textit{Ways to understand God's ``design'' in making only men priesthood bearers.}}
		\end{table}
		
		We'll start with sense \#1.
		This is the very strongest understanding of God's ``design'' in giving only men the priesthood---%
			according to this one,
				God just \textit{wanted} to,
					and we may not get further explanation.
		The difference between ``for either no reason at all'' and ``for reasons we can never know or understand'' is important,
			because the latter implies that there \textit{are} reasons,
				even if we can't know them.
		We'll talk about that difference in a couple posts.
		The most important thing here is that sense \#1 is like the placement of the our friend's kitchen island anywhere in her house,
			except it's even less restricted than that.
		God \textit{could} have done anything,
			and almost by caprice chose to make men the bearers of the priesthood.
			
		Now sense \#2.
		It's similar to the first sense,
			but look at what has changed---%
				the interpretation now includes reference to other things outside of God,
					and we can know what they are.
		The sense of design here is still very strong,
			because God had many options,
				but there is a weak constraint acting on God's choice.%
					\footnote{The ``constraint'' is false, which is one among several reasons to reject this interpretation.}
		This is more analogous to the location of the sink and oven constraining our friend in her placement of the island---%
			she still had many options,
				but fewer than before.
				
		Sense \#3 introduces the first real physical constraint on this list,
			and as we know from our example of the boulder and the semicircle,
				physical constraints suggest a weaker sense of design.
		On this understanding,
			God looks at the structure of eternity,
				sees some kind of constraint,
					and makes a choice on that basis.
		The physical constraint determines the design choice;
			for some reason,
				it couldn't have been otherwise.%
					\footnote{If this were true, you would wonder why the structure of eternity was like that; what could the reason be? If we don't have an answer, we're in a position similar to that of the first sense, where God is acting on unknowable reasons.}
		This situation has a counterpart with our friend's house,
			in constructing the kitchen with a semicircle to accommodate the physical constraint of the boulder.
		The physical constraint makes it a weak sense of design.
			
		talk about sense \#1, and then connect/transition to the rest of the paragraph below, then move on
	
		an extremely important point here is to make sure we distinguish what we mean by ``design'' when talking about God---since God bestowed the priesthood on JS, of course he ``designed'' or ``intended,'' in a certain sense, that only men receive it. But when we're talking about the teleological fallacy we are talking about a stronger sense of ``design''---it's more like, the sense we mean is almost as though God had a choice among several options, and he chose the one he did for certain reasons---either because he just wanted to, or in this case, because it might reflect something about the structure of the plan of salvation. If he were NOT bound by that structure, he could STILL be bound by OTHER constraints, such as what human civilization is like and what people are like, and could therefore "design" (in the weak sense) that only men receive it. But these are very different reasons behind the design---the constraint comes from a totally different place and they are of different strengths and different durations.
		
		design gets weaker -> indicates that you are dealing less and less with an aspect of eternity, with the structure of eternity
		
		connect to teleological fallacy?
			the likelihood of the teleological and structural fallacies increases as the sense of design becomes weaker; as we go down the rows of this table, we are more likely to commit those fallacies.
		
	-----
	\subsubsection{Varieties of constraint}
	-----
	
		"physical constraints" doesn't just include stuff in physics, it's human psychology, how civilizations work, everything like that.
			Very important to have this broad understanding of what a physical constraint is.
			note that the stuff about human psychology and civilizations and things have to be UNAVOIDABLE FACTS ABOUT HOW THINGS WORK in order to count as physical constraints; if it's local to one person or something, then it doesn't count as a physical constraint.
		but we need something like a rank ordering of how strong/lasting the constraints are; how much weight they carry in eternity, you might say.
			Any design that involves physical constraints is by definition weak.
			And you then rank order the designs within the physical constraints by how strong/lasting the constraints are:
				if the constraints are MORE PHYSICAL or MORE BASIC, such as the laws of physics or chemistry, then you have stronger design that we still classify as weak.
				If the constraints are LESS PHYSICAL, LESS BASIC, or have more to do with the facts of biology and psychology, then you get into the weakest of the weak designs.
					It's even weaker, maybe the weakest, when you are talking about the characteristics of a particular time, place, culture, people, group, or something like that.
					There's a rank order within designs that have to do with people too.
					In the table above, we reject sense \#2 because what it says about men compared to women is just false.
	
	-----
	\subsubsection{Varieties of constraint in the gospel}
	-----
	
		you might think that the structure helps us understand the reason why god created something a certain way.
			so if there is some cosmic necessity behind something, then that's a necessary structural part of the plan of salvation, and so God designed it that way because he really didn't have a choice.
			this is a very tempting way of thinking and also a significant source of error, because we usually don't distinguish clearly enough the grand, cosmic necessities (whatever they are) from the contingent necessities (if you will) that arise in this implementation of the plan of salvation. The former are structural aspects of the plan or universe; the latter are not, but do force God to set up the plan a certain way in this (or any) implementation of the gospel.
			I will return to this topic in a later essay, as it is extremely important.
			See the paragraph in 0004, structural fallacy, about explaining how a structural inference could be good.
			Structural inferences fail when the constraints aren't deep or lasting
	
	-----
	\subsubsection{On the possibility of physical constraints for God}
	-----
	
		need to talk about the very possibility that God is subject to physical constraints in the first place---this is not a traditional view of Christianity and it is one of the most important innovations of the restoration.
			That it is even possible for God to be constrained in this way is a very important philosophical development.
			The other way of thinking about it is, God just wills certain things to be a certain way